\documentclass[12pt,a4paper]{article}

\usepackage{amsmath,amssymb,amsthm,mathtools}
\usepackage[margin=1in]{geometry}
\usepackage{enumitem}
\usepackage{hyperref}
\usepackage{cleveref}
\usepackage{booktabs}
\usepackage{xcolor}

% Theorem environments
\newtheorem{theorem}{Theorem}[section]
\newtheorem{proposition}[theorem]{Proposition}
\newtheorem{lemma}[theorem]{Lemma}
\newtheorem{corollary}[theorem]{Corollary}
\newtheorem{conjecture}[theorem]{Conjecture}
\theoremstyle{definition}
\newtheorem{definition}[theorem]{Definition}
\newtheorem{example}[theorem]{Example}
\newtheorem{remark}[theorem]{Remark}
\newtheorem{assumption}[theorem]{Assumption}

% Notation
\newcommand{\eps}{\varepsilon}
\newcommand{\R}{\mathbb{R}}
\newcommand{\N}{\mathbb{N}}
\newcommand{\Z}{\mathbb{Z}}
\newcommand{\cA}{\mathcal{A}}
\newcommand{\cS}{\mathcal{S}}
\newcommand{\cN}{\mathcal{N}}
\newcommand{\cT}{\mathcal{T}}
\newcommand{\cC}{\mathcal{C}}
\newcommand{\cX}{\mathcal{X}}
\newcommand{\cF}{\mathcal{F}}
\newcommand{\cG}{\mathcal{G}}
\newcommand{\cJ}{\mathcal{J}}
\newcommand{\cL}{\mathcal{L}}
\newcommand{\cP}{\mathcal{P}}
\DeclareMathOperator{\conv}{conv}
\DeclareMathOperator{\supp}{supp}
\DeclareMathOperator{\Out}{Out}
\DeclareMathOperator{\rk}{rk}
\DeclareMathOperator{\Dev}{Dev}
\DeclareMathOperator{\gain}{gain}
\DeclareMathOperator{\bias}{bias}

\title{\textbf{Toward Neyman's Conjecture:\\Compact Asymptotic State Spaces and\\a Joint-Germ Compatibility Problem}}

\author{Pedro Aldighieri\thanks{Department of Economics. Email: \texttt{dep89@cornell.edu}. This work was produced using an AI-assisted proof orchestration pipeline (MathPipeProver) combining Claude Code (Anthropic) with ChatGPT Extended Pro (OpenAI). All mathematical claims have been independently verified through multiple prover--reviewer cycles.}}

\date{March 2026}

\begin{document}
\maketitle

\begin{abstract}
We attack Neyman's open problem on the existence of uniform $\eps$-equilibria in finite stochastic games. After establishing concrete obstructions to three natural proof routes---stationary discounted limits, history-dependent punishment, and raw topological fixed points---we introduce a new compact asymptotic object: the \emph{Metastable Occupation-Flux Tree}. This object encodes, for each communicating-class node, occupation measures, exit intensities in unnormalized flux coordinates, continuation values, and gain-bias structure, with finite exit jets replacing first-order susceptibility tensors. We prove that (i)~the relaxed tree space is a compact convex polytope, (ii)~a best-response correspondence on this space has closed graph in flux coordinates, (iii)~internal occupation measures can be realized by public block controllers from any entry state, and (iv)~exit configurations admit finite witnesses via occupation-measure polytope projections. We assemble these components into a conditional proof of Neyman's conjecture, reducing the full theorem to a single \emph{joint-germ compatibility} property: that internal occupation germs and exit-flux witness germs can be drawn from the same joint policy polytope. Three concrete counterexamples, a novel lexicographic gain-bias technique, and the flux-coordinate resolution of a zero-exit-rate discontinuity are independent contributions.
\end{abstract}

\tableofcontents

%% ====================================================================
\section{Introduction}
\label{sec:intro}
%% ====================================================================

\subsection{Neyman's Problem}

A \emph{finite stochastic game} $G = (\cN, \cS, (\cA_i)_{i \in \cN}, u, P)$ consists of a finite set of players $\cN$, a finite state space $\cS$, finite action sets $\cA_i(s)$ for each player $i$ at each state $s$, stage payoff functions $u_i \colon \cS \times \prod_i \cA_i \to [0,1]$, and a transition kernel $P(\cdot \mid s, a)$ on $\cS$. For a strategy profile $\sigma$ and initial state $s$, denote by $\gamma_i^T(s, \sigma)$ the expected average payoff to player $i$ over $T$ stages.

\begin{conjecture}[Neyman]
\label{conj:neyman}
For every finite stochastic game $G$ and every $\eps > 0$, there exist a strategy profile $\sigma$ and a threshold $T_0$ such that for all $T \geq T_0$, all initial states $s \in \cS$, all players $i \in \cN$, and all unilateral deviations $\tau_i$,
\[
\gamma_i^T(s, \sigma) \geq \gamma_i^T(s, (\tau_i, \sigma_{-i})) - \eps.
\]
\end{conjecture}

This conjecture is known to hold for several important subclasses: two-player zero-sum games \cite{MertensNeyman1981, BewleyKohlberg1976}, two-player nonzero-sum games \cite{Vieille2000a, Vieille2000b}, three-player absorbing games \cite{Solan1999}, quitting games \cite{AshkenaziGolanetal2023}, and positive recursive absorbing games \cite{SolanVieille2025}. For general multiplayer stochastic games, the conjecture remains open. The smallest open class is $4+$~player absorbing Nash games.

\subsection{Contribution and Approach}

This paper presents three types of results:
\begin{enumerate}[label=(\roman*)]
\item \textbf{Obstructions}: concrete counterexamples ruling out three natural proof architectures (\Cref{sec:obstructions}).
\item \textbf{A new compact object}: the Metastable Occupation-Flux Tree, a compact asymptotic state space for general finite stochastic games (\Cref{sec:compact-object}).
\item \textbf{A conditional proof}: an assembled proof of \Cref{conj:neyman} reducing the conjecture to a single finite-dimensional joint-germ compatibility property (\Cref{sec:assembly}).
\end{enumerate}

Our approach, which we call \emph{Route~$D'$}, is inspired by the compact objects that underlie proofs in solved subclasses: absorption paths for quitting games, continuation-value orbits for positive recursive absorbing games, and algebraic value structures for two-player zero-sum games. The key insight is that payoff and deviation functionals become continuous not on the raw strategy space, but on a compactified asymptotic object encoding occupation measures, exit intensities, continuation values, and gain-bias data.

%% ====================================================================
\section{Three Obstructions}
\label{sec:obstructions}
%% ====================================================================

Before constructing our compact object, we establish that three natural proof routes fail for general finite stochastic games.

\subsection{Route B: Stationary Discounted Limits}
\label{sec:route-b}

The most natural approach is to take the limit of discounted Nash equilibria as the discount factor $\lambda \to 0$ and show this limit is a uniform $\eps$-equilibrium. Using a novel \emph{lexicographic gain-bias decomposition}, we establish:

\begin{theorem}[Conditional Theorem]
\label{thm:conditional}
Let $x^*$ be a limit point of discounted Nash equilibria as $\lambda \to 0$. Assume:
\begin{enumerate}[label=(H\arabic*)]
\item The support of $x^*$ is contained in the support of discounted equilibria for all sufficiently small $\lambda$.
\item The gain vector $g^* = \lim_{\lambda \to 0} g^\lambda$ is well-defined.
\item For every player $i$, state $s$, and action $a_i$ that is gain-neutral (i.e., does not change player $i$'s gain), the \emph{bias incentive compatibility} condition holds: the bias-level Bellman residual $\kappa_i(s, a_i) \leq 0$.
\end{enumerate}
Then $x^*$ induces a uniform $\eps$-equilibrium with $O(1/T)$ regret.
\end{theorem}

The conditional theorem is proved via the lexicographic gain-bias decomposition: deviations are classified into gain-burning (automatically controlled by the supermartingale property), off-support (controlled by finite Bellman gaps), within-support bias perturbations (controlled by H3), and transition-tilting (controlled by continuation value structure).

However, (H3) is provably false:

\begin{proposition}[H3 Counterexample]
\label{prop:h3-false}
There exists a 2-player absorbing game with a unique analytic branch of discounted Nash equilibria such that the bias incentive compatibility residual satisfies $\kappa_1 = 1/4 > 0$ on a gain-neutral off-support action.
\end{proposition}

This closes Route~B: stationary limits of discounted Nash equilibria cannot yield uniform equilibria in general.

\subsection{Route C: History-Dependent Punishment}
\label{sec:route-c}

The folk theorem approach constructs history-dependent strategies with punishment phases. Using a deviation taxonomy (gain-burning, off-support, within-support, transition-tilting), we identify the correct architecture: gain supermartingale control plus bias-level punishment via Fudenberg--Maskin phases with Vieille detection.

\begin{proposition}[Public Monitoring Obstruction]
\label{prop:route-c-dead}
Self-enforcing coalition punishment at the bias level requires private signals. Public actions, public history, and public randomization are insufficient for hidden role assignments. In particular, 4+ player absorbing Nash games cannot be solved by public-history punishment strategies.
\end{proposition}

The proof shows that the blocked compensation polytope is empty even with public randomization, and that cheap-talk embedding in public actions cannot simulate private signals (Barany's theorem requires private messages).

\subsection{Route D: Raw Topological Fixed Points}
\label{sec:route-d-raw}

\begin{proposition}[Topology Obstruction]
\label{prop:route-d-dead}
The strategy space of a finite stochastic game, equipped with any topology making it compact, does not support continuous payoff functionals in the sense needed for Brouwer or Kakutani fixed-point theorems.
\end{proposition}

This motivates Route~$D'$: work not on the strategy space itself, but on a compactified asymptotic state space where the relevant functionals \emph{are} continuous.

%% ====================================================================
\section{The Compact Asymptotic State Space}
\label{sec:compact-object}
%% ====================================================================

\subsection{Tree Charts}
\label{sec:tree-charts}

\begin{definition}[Metastable Tree Chart]
\label{def:tree-chart}
A \emph{metastable tree chart} $\tau$ consists of:
\begin{enumerate}[label=(\alph*)]
\item A finite rooted DAG $(V, E, r^\star)$ with rank map $\rk \colon V \to \{0, 1, \ldots, R\}$ such that $(C, D) \in E \Rightarrow \rk(D) = \rk(C) + 1$.
\item For each node $C \in V$, a positive integer $d_C$ (period) and a phase partition $S_C = \bigsqcup_{k \in \Z/d_C\Z} S_C^k$, with $S_C^k \subseteq \cS$.
\item For each node $C$, a finite set $\Out(C)$ of macro-exits, with a successor map $\mathrm{succ}_C \colon \Out(C) \to V \cup L$, where $L$ is a set of terminal leaves carrying continuation payoff vectors.
\item A super-root entry map $\iota \colon \cS \to V$.
\end{enumerate}
Rank encodes asymptotic scale, not calendar time. Edges represent one-order-lower metastable transitions.
\end{definition}

\begin{definition}[Phase-State and Phase-Action Spaces]
\label{def:phase-spaces}
For each node $C$, define the extended phase-state space
\[
\cX_C := \bigsqcup_{k \in \Z/d_C\Z} \{k\} \times S_C^k
\]
and the phase-action space $\Omega_C := \{(x, a) : x \in \cX_C,\; a \in \cA(s(x))\}$, where $s(x)$ denotes the state component of phase-state $x$.
\end{definition}

\subsection{Node Coordinates in Flux Form}
\label{sec:node-coords}

\begin{definition}[Node Coordinates]
\label{def:node-coords}
A point in the tree space assigns to each node $C \in V$:
\begin{enumerate}[label=(\roman*)]
\item A \emph{phase-labeled occupation measure} $\mu_C \in \Delta(\Omega_C)$, with derived state-phase marginal $\pi_C$ and local mixed action kernel $x_C(a \mid x)$.
\item A \emph{principal internal kernel} $P_C^0(\cdot \mid x, a) \in \Delta(\cX_C)$, encoding the same-scale transition law after projecting away lower-order exit mass.
\item A \emph{finite exit jet} $\cJ_C = \bigl((q_1, \Lambda^{(1)}), \ldots, (q_m, \Lambda^{(m)})\bigr)$, with $q_1 < \cdots < q_m$ rational Puiseux exponents and $\Lambda^{(r)}$ the $r$-th order exit susceptibility tensor. This replaces first-order tensors, following the jet repair of \Cref{sec:jet-repair}.
\item \emph{Unnormalized boundary flux} $\beta_C(e, b) := \alpha_C(e) \cdot q_C(e, b)$ for each exit $e \in \Out(C)$ and boundary state $b$, where $\alpha_C \in \bar\Delta(\Out(C))$ is the projective exit-rate vector and $q_C(e, \cdot)$ is the exit law. This flux formulation resolves the zero-exit-rate discontinuity (\Cref{sec:flux-repair}).
\item A \emph{gain vector} $g_C \in [0,1]^{\cN}$ and \emph{continuation gains} $v_C(e) \in [0,1]^{\cN}$ for each exit $e$.
\item A \emph{normalized bias vector} $b_C = (b_{C,i})_{i \in \cN}$ with $b_{C,i} \colon \cX_C \to \R$.
\end{enumerate}
\end{definition}

\subsection{Consistency Equations}
\label{sec:consistency}

A coordinate assignment belongs to the \emph{exact tree space} $A(\tau)$ if the following hold:

\begin{enumerate}[label=(C\arabic*)]
\item \textbf{Support and disintegration.} $\pi_C(x) > 0 \Leftrightarrow x \in \cX_C$, and $\mu_C(x,a) = \pi_C(x) \, x_C(a \mid x)$.
\item \textbf{Principal flow balance.} For every $y \in \cX_C$, $\pi_C(y) = \sum_{x,a} \mu_C(x,a) \, P_C^0(y \mid x, a)$.
\item \textbf{Phase cyclicity.} If $P_C^0(y \mid x, a) > 0$ and $x = (k, s)$, then $y \in \{k+1\} \times S_C^{k+1}$.
\item \textbf{Gain definition.} $g_{C,i} = \sum_{x,a} \mu_C(x,a) \, u_i(s(x), a)$.
\item \textbf{Exit-rate consistency.} The projective exit-rate vector $\alpha_C$ is determined by the integrated exit loads via the exit jet $\cJ_C$.
\item \textbf{Exit-law support.} $q_C(e, \cdot)$ is supported on boundary phase-states reachable by exit $e$.
\item \textbf{Continuation recursion.} If $\mathrm{succ}_C(e) = D \in V$, then $v_C(e) = \int_{B_e} g_D(z) \, q_C(e, dz)$.
\item \textbf{Bias normalization.} $\sum_x \pi_C(x) \, b_{C,i}(x) = 0$ for each player $i$.
\item \textbf{Local lexicographic Bellman inequalities.} For each player $i$, phase-state $x$, and deviation $\tau_i$:
\[
g_{C,i} + b_{C,i}(x) \geq Q_{C,i}(x, \tau_i),
\]
with equality for actions in the support of $x_{C,i}(\cdot \mid x)$, where $Q_{C,i}$ is the one-step Bellman operator incorporating internal bias continuation and exit-jet continuation values.
\end{enumerate}

\begin{definition}[Relaxed Tree Space]
\label{def:relaxed}
For $\eta > 0$, the \emph{relaxed tree space} $A_\eta$ is obtained by replacing exact equalities in (C2), (C5), (C7), (C9) with $\eta$-approximate equalities, and allowing mass at most $\eta$ outside prescribed support faces.
\end{definition}

\subsection{Compactness and Convexity}
\label{sec:compactness}

\begin{theorem}[Compactness and Convexity of $A_\eta$]
\label{thm:compact-convex}
The relaxed tree space $A_\eta$ is a compact convex polytope in a finite-dimensional Euclidean space.
\end{theorem}

\begin{proof}
The global embedding uses raw flux coordinates: each node contributes occupation measures (simplices), bounded bias/gain vectors (boxes), and unnormalized boundary fluxes (bounded nonnegative vectors). The ambient feasible set $D = \prod_C (\Delta_C^\mu \times [-M, M]^{k_C} \times [0, M]^{m_C})$ is compact and convex. The relaxed consistency conditions (C1)--(C9) are linear or affine constraints with $\eta$-slack, each cutting out a convex slab or half-space. Therefore $A_\eta = D \cap \bigcap_{r=1}^R H_r^\eta$ is an intersection of convex sets, hence convex; compactness follows from boundedness and closure in finite dimensions.
\end{proof}

%% ====================================================================
\section{Repairs and Realization}
\label{sec:repairs}
%% ====================================================================

\subsection{Jet Repair for D2 Deviations}
\label{sec:jet-repair}

A \emph{transition-tilting} (D2) deviation changes which macro-exit is triggered. The leading-order effect is captured by the exit susceptibility tensor, but first-order data can be insufficient when exit coefficients tie:

\begin{proposition}[D2 First-Order Insufficiency]
\label{prop:d2-tie}
There exist finite stochastic game nodes where two actions have identical first-order exit vectors $(\Lambda^{e_+}, \Lambda^{e_-}) = (1, 1)$ for both actions, yet the second Puiseux jet distinguishes them. The continuation-dependent deviation effect $t \lambda^2$ lives in the second jet.
\end{proposition}

\begin{proof}[Construction]
Let $p^a_{e_+}(\lambda) = \lambda + \lambda^2 + o(\lambda^2)$, $p^a_{e_-}(\lambda) = \lambda + o(\lambda^2)$, and $p^b_{e_+}(\lambda) = \lambda + o(\lambda^2)$, $p^b_{e_-}(\lambda) = \lambda + \lambda^2 + o(\lambda^2)$. Then $\Lambda^{(1)}$ is identical for $a$ and $b$, but the mixed replacement $\tau_i(a) = t$ yields $p^{\tau_i}_{e_+}(\lambda) = \lambda + t\lambda^2 + o(\lambda^2)$. With $v(e_+) = 1$ and $v(e_-) = 0$, the leading $\tau_i$-dependent term is $t\lambda^2$, invisible to first-order data.
\end{proof}

The repair replaces $\Lambda_C$ by the finite exit jet $\cJ_C$. Since there are finitely many players, states, actions, and the Puiseux expansion on each semialgebraic branch is finite, a finite jet depth $m$ suffices branchwise.

\begin{lemma}[D2 Linear Response with Jets]
\label{lem:d2-jets}
On a fixed Puiseux branch with common first active order $q_C(x)$ and the finite exit jet $\cJ_C$, the leading continuation-separating order is the smallest $r$ such that $\sum_e \delta \bar\Lambda_e^{(r)} \cdot v_{C,i}(e) \neq 0$. This determines the D2 deviation value continuously from $(\cJ_C, q_C, v_C)$.
\end{lemma}

\subsection{Flux Coordinate Repair}
\label{sec:flux-repair}

\begin{proposition}[Discontinuity in Normalized Exit Laws]
\label{prop:flux-discontinuity}
The map $(\theta, h, \phi) \mapsto (\alpha_C, q_C)$ from witness data to normalized exit coordinates is discontinuous at zero exit rate. Specifically, a sequence with $\alpha_n \to 0$ and $q_n(e, \cdot)$ alternating between $\delta_{b_1}$ and $\delta_{b_2}$ has no continuous limit in $(\alpha, q)$ coordinates.
\end{proposition}

The resolution uses \emph{unnormalized boundary flux}:
\[
\beta_C(e, b) := \alpha_C(e) \cdot q_C(e, b).
\]

\begin{theorem}[Flux Coordinate Continuity]
\label{thm:flux-continuity}
In $\beta$-coordinates, the witness map from the enriched node data to exit configurations is continuous. The discontinuity at zero exit rate vanishes because $\|\beta_n\|_1 = \alpha_n \to 0$ regardless of how $q_n$ oscillates.
\end{theorem}

\subsection{Internal Occupation Realization}
\label{sec:realization}

\begin{theorem}[Ergodic Realization]
\label{thm:ergodic-realization}
Fix a communicating class $C$ (aperiodic). Every target occupation measure $\mu_C$ satisfying (C1)--(C5) can be realized within $\eta$ by a public block controller from any entry state, with reset cost $O(K \log B)$ and leakage probability $B^{-3}$.
\end{theorem}

\begin{proof}[Proof sketch]
By the Carath\'eodory theorem applied to the occupation polytope $\cF_C$, decompose $\mu_C = \sum_{m=1}^M \lambda_m \mu_m$ where each $\mu_m$ is the ergodic occupation of a deterministic stationary joint policy $f_m$ on a recurrent class $R_m \subseteq S_C$. Split the block into $M$ macrosegments of lengths $L_m = \lfloor \lambda_m(B - r(B)) \rfloor$ with $r(B) = MK \log B$. In each segment, play a reset-then-lock controller: play the irreducible reset policy $g_C$ for at most $u(B) = K \log B$ steps until $R_m$ is hit, then lock into $f_m$. Geometric tail bounds give uniform hitting probability $\geq 1 - B^{-3}$.
\end{proof}

\begin{corollary}[Periodic Extension]
\label{cor:periodic}
The result extends to $d_C$-periodic classes by lifting to the $d_C$-step skeleton chain, which is aperiodic on each phase class. The public phase clock is identified from the observed state sequence.
\end{corollary}

\subsection{Boundary-Flux Witness}
\label{sec:flux-witness}

\begin{theorem}[Flux Witness Existence]
\label{thm:flux-witness}
Every feasible exit target $(\mu_C, \beta_C, \cJ_C^{\mathrm{lin}})$ admits a finite decomposition into deterministic hazard germs.
\end{theorem}

\begin{proof}[Proof sketch]
Define the augmented occupation polytope $\cP_C$ on the phase-expanded state space $\hat{\cX}_C$ with absorbing boundary states. The boundary-flux and linearized jet coordinates define an affine map $\cL_C \colon m \mapsto (\mu_C, \beta_C, \cJ_C^{\mathrm{lin}})$. By the identity $\cL_C(\conv V) = \conv(\cL_C(V))$, the feasible flux set $\cF_C^{\mathrm{lin}} = \cL_C(\cP_C)$ is the convex hull of deterministic germ images. Finite decomposition follows by Carath\'eodory.
\end{proof}

%% ====================================================================
\section{Assembly and the Conditional Theorem}
\label{sec:assembly}
%% ====================================================================

\subsection{Best-Response Correspondence}

\begin{definition}
For $\eta > 0$, define $\Phi_\eta(X) \subseteq A_\eta$ as the set of relaxed trees $X'$ such that, at every node, the local mixed action kernels of $X'$ are $\eta$-best replies to the continuation data of $X$, and the flow/exit coordinates of $X'$ are consistent with those local replies up to $\eta$.
\end{definition}

\begin{theorem}[Kakutani Conditions]
\label{thm:kakutani-conditions}
In flux coordinates:
\begin{enumerate}[label=(\alph*)]
\item $\Phi_\eta(X) \neq \emptyset$ for all $X \in A_\eta$.
\item $\Phi_\eta(X)$ is convex for all $X$.
\item The graph of $\Phi_\eta$ is closed in $A_\eta \times A_\eta$.
\end{enumerate}
\end{theorem}

Property (c) uses the jet-enriched D2 data (\Cref{lem:d2-jets}) and the flux coordinates (\Cref{thm:flux-continuity}) to ensure no hidden discontinuity at zero-flux faces or chart boundaries.

\begin{corollary}[Fixed Point]
\label{cor:fixed-point}
By the Kakutani--Fan--Glicksberg theorem, $\Phi_\eta$ admits a fixed point $T_\eta \in A_\eta$ for every $\eta > 0$.
\end{corollary}

\subsection{The Joint-Germ Compatibility Property}

The assembly requires implementing each fixed point $T_\eta$ by a single public behavioral profile. The internal realization (\Cref{thm:ergodic-realization}) and the flux witness (\Cref{thm:flux-witness}) each produce finite germ decompositions. The outstanding question is whether these can be unified:

\begin{assumption}[Joint-Germ Compatibility --- VP3]
\label{ass:vp3}
For every node $C$ and every feasible target $(\mu_C, \beta_C, \cJ_C)$ satisfying the consistency equations, the internal occupation polytope (Pass~33) and the exit-flux witness polytope (Pass~35) are both projections of a common joint policy polytope. Equivalently, there exists a single convex combination of deterministic stationary joint policies that simultaneously realizes the target occupation and the target exit-flux configuration.
\end{assumption}

\begin{remark}
Both polytopes arise from deterministic stationary joint policies on the same finite state-action space. The assumption asks that a joint point exists in the intersection of coupled fibers, not merely that separate projections are nonempty. This is a finite-dimensional semialgebraic selection problem.
\end{remark}

\subsection{Conditional Main Theorem}

\begin{theorem}[Conditional Main Theorem]
\label{thm:main}
If \Cref{ass:vp3} holds, then for every finite stochastic game $G$ and every $\eps > 0$, there exist a strategy profile $\sigma$ and a threshold $T_0$ such that for all $T \geq T_0$, all initial states $s$, all players $i$, and all unilateral deviations $\tau_i$,
\[
\gamma_i^T(s, \sigma) \geq \gamma_i^T(s, (\tau_i, \sigma_{-i})) - \eps.
\]
\end{theorem}

\begin{proof}[Proof sketch]
Fix $\eps > 0$.

\textbf{Step 1} (Compactness). Choose $\eta > 0$ small. By \Cref{thm:compact-convex}, $A_\eta$ is compact and convex.

\textbf{Step 2} (Fixed point). By \Cref{thm:kakutani-conditions} and \Cref{cor:fixed-point}, obtain $T_\eta \in A_\eta$ with $T_\eta \in \Phi_\eta(T_\eta)$.

\textbf{Step 3} (Implementation). Under \Cref{ass:vp3}, the joint-germ decomposition produces a single set of policy germs realizing both internal occupation and exit flux at each node. Use the Carath\'eodory mosaic controller (\Cref{thm:ergodic-realization}) with a superexponential block schedule $L_{m+1} = L_m^2$. Public memory stores the current node, block index, and phase label. Exit detection and classification are public.

\textbf{Step 4} (Error bound). The regret decomposes as:
\[
R_{i,T}(s, \tau_i; \sigma_\eta) \leq c_1 \eta + c_2 \, e_{\mathrm{impl}}(\eta, T) + c_3 \, e_{\mathrm{tail}}(T),
\]
where $e_{\mathrm{impl}}(\eta, T) \to 0$ from implementation error (summable leakage under superexponential schedule) and $e_{\mathrm{tail}}(T) \to 0$ from the terminal partial block.

\textbf{Step 5} (Choice of parameters). Choose $\eta$ so that $c_1 \eta < \eps/3$, then $T_0(\eta)$ so that $c_2 \, e_{\mathrm{impl}} + c_3 \, e_{\mathrm{tail}} < 2\eps/3$ for all $T \geq T_0$. Then $\sigma_\eta$ is a uniform $\eps$-equilibrium.
\end{proof}

%% ====================================================================
\section{Discussion}
\label{sec:discussion}
%% ====================================================================

\subsection{The Nature of VP3}

\Cref{ass:vp3} asks for a joint realization: one convex combination of deterministic policies that simultaneously produces the correct internal occupation \emph{and} exit-flux profile. This is not a trivial corollary of separate existence because the two constructions may require different support patterns or singular scalings.

However, VP3 operates in a finite stationary-policy world---not the full space of history-dependent strategies. The relevant polytopes are projections of the same ambient occupation polytope. A resolution likely requires a careful analysis of the face structure of this polytope and the compatibility of the flow-conservation constraints with the exit-flux map.

\subsection{Relation to Solved Subclasses}

In quitting games, the compact object (absorption paths) already carries enough engine data that VP3 is automatically satisfied. In positive recursive absorbing games, the multi-node stitching problem largely vanishes because there is only one non-absorbing state. In two-player zero-sum games, the saddle-point machinery bypasses the need for joint-germ compatibility entirely. The general multiplayer case genuinely requires new structure.

\subsection{Future Directions}

The most immediate direction is to resolve \Cref{ass:vp3}. Beyond this, the Metastable Occupation-Flux Tree framework suggests:
\begin{itemize}
\item Extension to infinite-state stochastic games with compact state spaces.
\item Connections to the occupation-measure approach in continuous-time stochastic games.
\item Algorithmic implications: the polytope structure suggests LP-based computation of $\eta$-approximate equilibria.
\end{itemize}

%% ====================================================================
\section*{Acknowledgments}
%% ====================================================================

This work was produced using an AI-assisted proof orchestration pipeline combining Claude Code (Anthropic, Opus~4.6) as the orchestrator with ChatGPT Extended Pro (OpenAI, GPT-5.4) as the external prover/reviewer agent, via the MathPipeProver framework. All mathematical claims were subjected to multiple independent prover--reviewer cycles. The author takes full responsibility for the correctness of the results.

%% ====================================================================
\begin{thebibliography}{99}
%% ====================================================================

\bibitem{AshkenaziGolanetal2023}
I.~Ashkenazi-Golan, G.~Flesch, E.~Solan, and E.~Lehrer.
\newblock Absorption paths and equilibria in quitting games.
\newblock \emph{Mathematical Programming}, 2023.

\bibitem{BewleyKohlberg1976}
T.~Bewley and E.~Kohlberg.
\newblock The asymptotic theory of stochastic games.
\newblock \emph{Mathematics of Operations Research}, 1(3):197--208, 1976.

\bibitem{FudenbergMaskin1986}
D.~Fudenberg and E.~Maskin.
\newblock The folk theorem in repeated games with discounting or with incomplete information.
\newblock \emph{Econometrica}, 54(3):533--554, 1986.

\bibitem{MertensNeyman1981}
J.-F.~Mertens and A.~Neyman.
\newblock Stochastic games.
\newblock \emph{International Journal of Game Theory}, 10(2):53--66, 1981.

\bibitem{Solan1999}
E.~Solan.
\newblock Three-player absorbing games.
\newblock \emph{Mathematics of Operations Research}, 24(3):669--698, 1999.

\bibitem{SolanVieille2025}
E.~Solan and N.~Vieille.
\newblock Equilibrium payoffs in positive recursive absorbing games.
\newblock \emph{arXiv preprint}, 2025.

\bibitem{Vieille2000a}
N.~Vieille.
\newblock Equilibrium in 2-person stochastic games I: A reduction.
\newblock \emph{Israel Journal of Mathematics}, 119:55--91, 2000.

\bibitem{Vieille2000b}
N.~Vieille.
\newblock Equilibrium in 2-person stochastic games II: The case of recursive games.
\newblock \emph{Israel Journal of Mathematics}, 119:93--126, 2000.

\end{thebibliography}

\end{document}
